\documentclass[12pt,a4paper]{article}
\usepackage{ctex}
\usepackage{amsmath,amscd,amsbsy,amssymb,latexsym,url,bm,amsthm}
\usepackage{epsfig,graphicx,subfigure}
\usepackage{enumitem,balance}
\usepackage{wrapfig}
\usepackage{mathrsfs,euscript}
\usepackage[usenames]{xcolor}
\usepackage{hyperref}
\usepackage[vlined,ruled,commentsnumbered,linesnumbered]{algorithm2e}

\newtheorem{theorem}{Theorem}
\newtheorem{lemma}[theorem]{Lemma}
\newtheorem{proposition}[theorem]{Proposition}
\newtheorem{corollary}[theorem]{Corollary}
\newtheorem{exercise}{Exercise}
\newtheorem*{solution}{Solution}
\newtheorem{definition}{Definition}
\theoremstyle{definition}


%\numberwithin{equation}{section}
%\numberwithin{figure}{section}

\renewcommand{\thefootnote}{\fnsymbol{footnote}}

\newcommand{\postscript}[2]
 {\setlength{\epsfxsize}{#2\hsize}
  \centerline{\epsfbox{#1}}}

\renewcommand{\baselinestretch}{1.0}

\setlength{\oddsidemargin}{-0.365in}
\setlength{\evensidemargin}{-0.365in}
\setlength{\topmargin}{-0.3in}
\setlength{\headheight}{0in}
\setlength{\headsep}{0in}
\setlength{\textheight}{10.1in}
\setlength{\textwidth}{7in}
\makeatletter \renewenvironment{proof}[1][Proof] {\par\pushQED{\qed}\normalfont\topsep6\p@\@plus6\p@\relax\trivlist\item[\hskip\labelsep\bfseries#1\@addpunct{.}]\ignorespaces}{\popQED\endtrivlist\@endpefalse} \makeatother
\makeatletter
\renewenvironment{solution}[1][Solution] {\par\pushQED{\qed}\normalfont\topsep6\p@\@plus6\p@\relax\trivlist\item[\hskip\labelsep\bfseries#1\@addpunct{.}]\ignorespaces}{\popQED\endtrivlist\@endpefalse} \makeatother



\begin{document}
\noindent

%========================================================================
\noindent\framebox[\linewidth]{\shortstack[c]{
\Large{\textbf{Exercise Sheet 15}}\vspace{1mm}\\
Discrete Mathematics by Hongfei Fu, 2023.11.27}}

\begin{center}
	\footnotesize{\color{black} \quad Name:\_\_\_\_\_\_\_\_\_  \quad Student ID:\_\_\_\_\_\_\_\_\_ \quad Class: \_\_\_\_\_\_\_\_\_\_\_\_} \\


\end{center}

\begin{enumerate}


\item ([R], Page 705, Exercise 36) 
In Exercises 36 determine whether the given graph has a Hamilton circuit. If it does, find such a circuit. If it does not, give an argument to show why no such circuit exists.\\
\includegraphics*[scale=0.7]{36@p705.png}\\

\textbf{Answer Area:}
\vspace{6cm}


\item ([R], Page 705, Exercise 40) 
Does the graph in Exercise 33 have a Hamilton path? If so, find such a path. If it does not, give an argument to show why no such path exists.
\includegraphics*[scale=0.7]{33@p705.png}\\


\textbf{Answer Area:}
\vspace{6cm}

\item ([R], Page 706, Exercise 49, \textbf{2pt}) 
Show that there is a Gray code of order $n$ whenever $n$ is
a positive integer, or equivalently, show that the $n$-cube
$Q_{n}, n>1,$ always has a Hamilton circuit. [Hint: Use
mathematical induction. Show how to produce a Gray
code of order $n \text { from one of order } n-1 .]$

\textbf{Answer Area:}
\vspace{9cm}

\iffalse
\item ([R], Page 706, Exercise 55) 
Show that a bipartite graph with an odd number of vertices does not have a Hamilton circuit.

\textbf{Answer Area:}
\vspace{6cm}
\fi

\item(Exercise 48, \textbf{2pt})
Can you find a simple connected graph with $n$ vertices with $n\ge3$ that does not have a Hamilton circuit, yet the degree of every vertex in the graph is at least (n - 1)/2?

\textbf{Answer Area:}
\vspace{12cm}

\item(Exercise 46, \textbf{1pt})
Show that any subgraph of the Petersen graph below obtained by deleting a vertex has a Hamilton circuit. 
\includegraphics*[scale=0.7]{46.png}\\
\textbf{Answer Area:}
\vspace{8cm}

\item(Exercise 67, \textbf{2pt})
Show that this graph does not have a Hamilton circuit
\includegraphics*[scale=0.7]{67.png}\\
\textbf{Answer Area:}

\end{enumerate}
%========================================================================
\end{document}
